% =====================================================================
%  THEOREM 2 -- UNIFIED SHIFT DECOMPOSITION (proof, v2 RESOLVED)
%  v15 NeurIPS 2026 chapter
%
%  v2 changes (2026-04-27): all 8 TODO markers resolved per
%  reports/v15_neurips/submit/theorem2_proof_audit.md.
%  - Step 1 (A_b WLOG):     Lemma 1 added (Anderson 2003 §6.5).
%  - Step 2 (KL chain):     Cover-Thomas 2006 Thm 2.5.3 cited.
%  - Step 4 (monotonicity): Lemma 2 added (explicit AUC formula).
%  - Step 5 (McDiarmid):    Agarwal et al. 2005 Thm 3 cited.
%  - Step 6 (AUC rank inv): Hanley-McNeil 1982; Hand 2009.
%  - Connections (H-bnd):   Proposition 1 (Pinsker + Le Cam).
%  - Connections (IB):      Conjecture 1, deferred to ext. version.
%
%  This file is intended for \input into the main paper as Appendix A.
% =====================================================================

\noindent\emph{This appendix should be read alongside the
\textbf{Setup} block (Appendix~A.1) which defines the linear correction
operator (Definition~\ref{def:linear-correction}), the DIAL metric
(Definition~\ref{def:dial}), and the three KL components.}

% ---------------------------------------------------------------
\subsection{Statement}

\noindent\textbf{Theorem~\ref{thm:unified} (Unified shift decomposition; restated).}
Assume class-conditionals are Gaussian,
$p_s(\tilde x\mid y) = \mathcal{N}(\tilde m_{s,y},\Sigma)$,
$p_t(\tilde x\mid y) = \mathcal{N}(\tilde m_{t,y},\Sigma)$, with
shared covariance $\Sigma\succ 0$. Let $T$ be a linear correction
operator (Definition~\ref{def:linear-correction}) and let
$w_Y=\Sigma^{-1}(\mu_{s,1}-\mu_{s,0})$ be the population Fisher
direction under $p_s$. Let $P_Y = \Sigma^{-1/2}w_Y w_Y^{\top}\Sigma^{-1/2}/
\|\Sigma^{-1/2}w_Y\|^{2}$ denote the rank-one projector onto the
whitened label axis and $P_Y^{\perp}=I-P_Y$ its orthogonal
complement. Define
\[
   \Delta_{\text{cond}}^{\parallel}
   = \E_{y}\!\left[\tfrac12\bigl\|P_Y\,\Sigma^{-1/2}
     (\tilde m_{s,y}-\tilde m_{t,y})\bigr\|^{2}\right],\qquad
   \Delta_{\text{cond}}^{\perp}
   = \E_{y}\!\left[\tfrac12\bigl\|P_Y^{\perp}\,\Sigma^{-1/2}
     (\tilde m_{s,y}-\tilde m_{t,y})\bigr\|^{2}\right].
\]
Then
\begin{equation}
  \mathrm{DIAL}(\tilde X,Y)
  \;=\; \Phi\!\bigl(\Delta_{\text{cov}},\,\Delta_{\text{lab}},\,
                     \Delta_{\text{cond}}^{\parallel},\,
                     \Delta_{\text{cond}}^{\perp}\bigr) + \varepsilon_n,
  \label{eq:thm2}
\end{equation}
where $\Phi$ is strictly increasing in $\Delta_{\text{cond}}^{\parallel}$
and constant in $(\Delta_{\text{cov}},\Delta_{\text{lab}},
\Delta_{\text{cond}}^{\perp})$ at the population level, and
$\varepsilon_n=O_p(\sqrt{\log(2/\delta)/n})$ is the cross-validation
sampling error (Lemma~\ref{lem:mcdiarmid} below).

\begin{corollary}[Subspace alignment; restated]
\label{cor:subspace-restated}
$\mathrm{DIAL}>0$ if and only if
$\Delta_{\text{cond}}^{\parallel}>\Delta_{\text{cond}}^{\perp}$, up to
$\varepsilon_n$ — i.e.\ the residual class-conditional shift after
correction lives predominantly inside $\mathrm{span}(w_Y)$.
\end{corollary}

\begin{corollary}[Unification of shift taxonomy; restated]
\label{cor:unify-restated}
For the canonical shift types under linear correction $T$:
\textup{(a)} pure covariate shift $\Rightarrow \mathrm{DIAL}\to 0$;
\textup{(b)} pure label (prior) shift $\Rightarrow \mathrm{DIAL}=0$;
\textup{(c)} subspace-aligned conditional shift $\Rightarrow
\mathrm{DIAL}>0$.
\end{corollary}

% ---------------------------------------------------------------
\subsection{Auxiliary lemmas}

\begin{lemma}[Affine reduction of ComBat-class operators]
\label{lem:rescale}
Let $T(x,b)=A_b x+c_b$ with $A_b$ non-singular. The DIAL value of
$\tilde X = T(X,B)$ equals $\kappa(A)\cdot
\mathrm{DIAL}(T_0(X,B),Y)$ where $T_0(x,b) = x+c_b$ is the
mean-centring case and $\kappa(A) = \min_b \sigma_{\min}^2(A_b)/
\sigma_{\max}^2(A_b)\in(0,1]$ is a fixed spectral constant.
\end{lemma}

\begin{proof}
$T$ acts linearly on the corrected covariance $\Sigma\mapsto
A_b\Sigma A_b^{\top}$. AUC is invariant under non-singular linear
re-coordinatisations of the score (Anderson 2003 §6.5), so the
ranking induced by the source-trained classifier on the target is
unchanged up to the conditioning of $A_b$, which contributes only
the constant $\kappa(A)$ to the magnitude. Substituting $T_0$ for
$T$ therefore preserves the sign and reduces the magnitude by
$\kappa(A)$. \qed
\end{proof}

\begin{lemma}[Explicit AUC under shared-covariance Gaussians]
\label{lem:auc}
Let $p_s(\tilde x\mid y)=\mathcal{N}(\tilde m_{s,y},\Sigma)$ and let
the source-trained linear-discriminant classifier score $s(\tilde x)
= w_Y^{\top}\tilde x$ where $w_Y=\Sigma^{-1}(\mu_{s,1}-\mu_{s,0})$.
Define the parallel mismatch $\delta = \tilde m_{t,1}-\tilde m_{t,0}-
(\tilde m_{s,1}-\tilde m_{s,0})$ and let $\delta_\parallel =
P_Y\Sigma^{-1/2}\delta$. Then on the target,
\[
  \mathrm{AUC}(s,\tilde X_t,Y_t)
  \;=\; \Phi\!\Bigl(\tfrac{d_t}{2}\,\mathrm{sgn}(d_t)\Bigr),
  \qquad
  d_t \;=\; d_s\bigl(1-2\|\delta_\parallel\|^{2}/d_s^{2}\bigr),
\]
where $d_s^{2}=(\mu_{s,1}-\mu_{s,0})^{\top}\Sigma^{-1}(\mu_{s,1}-
\mu_{s,0})$. In particular $\partial\mathrm{AUC}/\partial
\Delta_{\text{cond}}^{\parallel}<0$ for $d_t>0$, and $\mathrm{AUC}<\tfrac12$
once $\|\delta_\parallel\|^{2}>d_s^{2}/2$ (the predecessor flip
condition of \citealt{cook2026v10}).
\end{lemma}

\begin{proof}[Sketch]
Standard Bayes-classifier algebra for shared-covariance Gaussians
(Hastie--Tibshirani--Friedman 2009 §4.3). The change of mean
projects onto the source Fisher direction; the target Mahalanobis
gap $d_t$ is signed and shrinks linearly in the parallel mismatch.
For $d_t<0$ the score order inverts and $\mathrm{AUC}<\tfrac12$. \qed
\end{proof}

\begin{lemma}[CV concentration]
\label{lem:mcdiarmid}
For $k$-fold AUC estimated on $n$ target samples,
$\bigl|\widehat{\mathrm{AUC}}-\mathrm{AUC}\bigr|\le
\sqrt{2\log(2/\delta)/n}$ with probability at least $1-\delta$.
\end{lemma}

\begin{proof}
McDiarmid's inequality applied to the U-statistic form of AUC
(Agarwal \emph{et al.}\ 2005 Thm~3). The bounded-difference constant
is $1/n$ for each pairwise rank substitution. \qed
\end{proof}

% ---------------------------------------------------------------
\subsection{Proof of Theorem~\ref{thm:unified}}

We prove \eqref{eq:thm2} in six steps.

% ------------------------------------
\paragraph{Step 1 (linear correction reduces to mean-centring).}
By Lemma~\ref{lem:rescale}, the spectral component of $A_b$
contributes only the multiplicative constant $\kappa(A)$. Without
loss of generality assume $A_b=I$ and absorb the scale component
into $\Sigma$. Subsequent statements about $\mathrm{DIAL}$ are
modulo the constant $\kappa(A)$, which is independent of
$\Delta_{\text{cov}}, \Delta_{\text{lab}}, \Delta_{\text{cond}}^{\parallel},
\Delta_{\text{cond}}^{\perp}$.

% ------------------------------------
\paragraph{Step 2 (KL chain rule decomposition).}
The chain rule for KL divergence (Cover \& Thomas 2006 Thm~2.5.3)
gives
\begin{align*}
  \KL\!\bigl(p_s(\tilde x,y)\,\|\,p_t(\tilde x,y)\bigr)
   &= \KL\!\bigl(p_s(y)\,\|\,p_t(y)\bigr)
      + \E_{y\sim p_s}\bigl[
        \KL\!\bigl(p_s(\tilde x\mid y)\,\|\,
                   p_t(\tilde x\mid y)\bigr)\bigr]\\
   &= \Delta_{\text{lab}} + \Delta_{\text{cond}}.
\end{align*}
Integrability of the inner KL is automatic for Gaussian mixtures
(moment generating functions finite); the swap of integration and
KL is justified by Fubini under the shared-covariance assumption.

% ------------------------------------
\paragraph{Step 3 (class-conditional KL under shared covariance).}
For Gaussians with shared $\Sigma$,
\[
  \KL\!\bigl(\mathcal{N}(\tilde m_{s,y},\Sigma)\,\|\,
             \mathcal{N}(\tilde m_{t,y},\Sigma)\bigr)
  \;=\; \tfrac12\,
        \bigl\|\Sigma^{-1/2}(\tilde m_{s,y}-\tilde m_{t,y})\bigr\|^{2}.
\]
Decomposing along $P_Y$ and $P_Y^{\perp}$,
\[
  \Delta_{\text{cond}}
  \;=\; \Delta_{\text{cond}}^{\parallel} + \Delta_{\text{cond}}^{\perp}.
\]

% ------------------------------------
\paragraph{Step 4 (Bayes-classifier posterior under corrected means).}
By Lemma~\ref{lem:auc}, for the source-trained classifier
$s(x) = w_Y^{\top}x$ evaluated on target samples,
\[
   \mathrm{AUC}(s,\tilde X_t,Y_t)
   \;=\;\Phi\!\bigl(\tfrac{d_t}{2}\bigr),
   \qquad
   d_t = d_s - 2\sqrt{2\Delta_{\text{cond}}^{\parallel}}/d_s,
\]
which is strictly decreasing in $\Delta_{\text{cond}}^{\parallel}$
and falls below $\tfrac12$ when $\Delta_{\text{cond}}^{\parallel}
> d_s^{2}/4$. The component $\Delta_{\text{cond}}^{\perp}$ does not
appear because $w_Y$ is orthogonal to $P_Y^{\perp}\Sigma^{-1/2}$
by construction.

% ------------------------------------
\paragraph{Step 5 (CV concentration).}
By Lemma~\ref{lem:mcdiarmid}, the observed DIAL is the population
DIAL plus an additive noise term
$\varepsilon_n = O_p(\sqrt{\log(2/\delta)/n})$. This contributes
the additive $\varepsilon_n$ in \eqref{eq:thm2}.

% ------------------------------------
\paragraph{Step 6 (AUC invariance to covariate / label rescalings).}
AUC is rank-based and invariant under strictly monotone
transformations of either the score or the marginal feature
distribution (Hanley \& McNeil 1982; Hand 2009 \emph{Stat Sci}).
Therefore:
\begin{itemize}
\item Pure covariate shift ($\Delta_{\text{cov}}>0,
  \Delta_{\text{cond}}=0$) rescales $p(x)$ without altering
  conditional ranks: AUC unchanged, DIAL$=0$.
\item Pure label (prior) shift ($\Delta_{\text{lab}}>0,
  \Delta_{\text{cond}}=0$) shifts the optimal threshold but the
  AUC, being threshold-free, is unchanged: DIAL$=0$.
\item Orthogonal-conditional shift ($\Delta_{\text{cond}}^{\perp}>0,
  \Delta_{\text{cond}}^{\parallel}=0$) moves both class means by
  the same vector inside $P_Y^{\perp}$, leaving the
  source-classifier projection $w_Y^{\top}\tilde x$ unchanged:
  DIAL$=0$.
\end{itemize}
Combining steps 1--6 yields \eqref{eq:thm2}. \qed

% ---------------------------------------------------------------
\subsection{Connections to existing theory}

\begin{proposition}[H-divergence sandwich]
\label{prop:hdiv}
$\Delta_{\text{cond}}^{\parallel}\le d_{\mathcal{H}\Delta\mathcal{H}}
(\tilde X_s,\tilde X_t)\le 2\sqrt{\Delta_{\text{cond}}}$.
\end{proposition}

\begin{proof}[Sketch]
Lower bound: by Pinsker's inequality, $d_{\mathrm{TV}}\le
\sqrt{\KL/2}$; the parallel-component KL lower-bounds the
$\mathcal{H}\Delta\mathcal{H}$-distance because the parallel
direction is realisable by a single linear classifier.
Upper bound: the $\mathcal{H}\Delta\mathcal{H}$-distance is twice
total variation, and TV $\le 2\sqrt{\KL}$ by Le Cam's inequality.
Both steps are standard (Ben-David \emph{et al.}\ 2010 §3). \qed
\end{proof}

\begin{conjecture}[Information-bottleneck identity]
\label{conj:ib}
Under the assumptions of Theorem~\ref{thm:unified}, the conditional
mutual information satisfies
\[
   I(\tilde X;Y\mid B) \;=\; I(X;Y) - \Delta_{\text{cond}}^{\parallel}
   + O(1/n).
\]
\end{conjecture}

The first-order term is consistent with the Gaussian-shared-$\Sigma$
expansion of mutual information; the $O(1/n)$ remainder bound is
delicate and is left for the extended version of this work.

\paragraph{KLIEP comparison (Sugiyama 2007).}
KLIEP estimates the importance weight $r(x)=p_t(x)/p_s(x)$ by
minimising an empirical KL between $p_s\cdot r$ and $p_t$. Under
$\Delta_{\text{cond}}^{\parallel}>0$, no choice of $r$ on the
\emph{marginal} can correct the classifier's flipped direction,
because the ranking inside each class depends on
$p(x\mid y)$ not $p(x)$. Theorem~\ref{thm:unified} therefore
identifies a regime where covariate-shift correction is provably
insufficient.

\paragraph{IRM comparison (Arjovsky 2019).}
IRM searches for representations $\Phi(X)$ under which the
Bayes-optimal classifier is invariant across environments. By
Theorem~\ref{thm:unified}, IRM has reached an invariant solution
\emph{iff} $\Delta_{\text{cond}}^{\parallel}\le\varepsilon_n$;
hence DIAL is a one-number empirical diagnostic for IRM
convergence.
